\documentclass{article}

\usepackage[T1]{fontenc}

\usepackage{booktabs,graphicx,siunitx}

\begin{document}

\section{Nested ANOVA}

\begin{tabular}{llll}
\toprule
Group & N & Mean & Std. Dev. \\
\midrule
north & 20 & 32.48 & 7.47 \\
south & 20 & 33.84 & 6.91 \\
west & 20 & 37.81 & 7.79 \\
\bottomrule
\end{tabular}

\begin{tabular}{lllllll}
\toprule
Source & SS & df & MS & F & p & $\omega$\textsuperscript{2} [95\% CI] \\
\midrule
school & 307.07 & 2 & 153.53 & 1.48 & .357 & 0.06 [0.00, 1.00] \\
classroom (nested in school) & 311.23 & 3 & 103.74 & 1.99 & .126 & 0.04 [0.00, 1.00] \\
\bottomrule
\end{tabular}

\includegraphics{figures/01_nested-anova/figure_grouped-means.png}

Used when one factor sits inside another (e.g. classes within schools). The top factor is tested against variation among its nested units, not raw residuals --- a significant F means the top-level groups differ beyond the nested-unit variability. Your significance threshold ($\alpha$) is 0.05.

A nested ANOVA for A gave F(2,3)=1.48, p = .357, $\omega$\textsuperscript{2}=.06 [.00, 1.00].

Statistical basis: Nested ANOVA (upper factor tested against the nested-unit mean square) (Fisher, R. A. (1925). Statistical Methods for Research Workers. Oliver \& Boyd.); $\omega$\textsuperscript{2} effect size (Ben-Shachar MS, Lüdecke D, Makowski D (2020). "effectsize: Estimation of Effect Size Indices and Standardized Parameters." Journal of Open Source Software, 5(56), 2815.).

\end{document}